% !TeX root = ../main.tex
% !TeX program = pdflatex
% !TeX encoding = UTF-8

\chapter{Mô hình hóa bài toán}
\label{chap:problem-modeling}

\section{Mô hình đồ thị}

Mạng đường được biểu diễn bằng đồ thị có hướng, có trọng số $G=(V,E)$. Tập
$V$ gồm các địa điểm và tập $E$ gồm các đoạn đường nối trực tiếp giữa chúng.
Đồ thị có hướng cho phép mô tả đường một chiều hoặc điều kiện di chuyển khác
nhau theo hai chiều. Nếu có thể đi theo cả hai chiều, dữ liệu lưu hai cạnh có
hướng riêng biệt.

Hình~\ref{fig:example-graph} minh họa mô hình trên một đồ thị nhỏ trước khi
dataset thực tế được trình bày ở Chương~\ref{chap:dataset}.

\begin{figure}[h!]
\centering
\begin{tikzpicture}[
  node distance=3cm,
  every node/.style={circle, draw, minimum size=1cm, font=\small},
  edge/.style={-{Latex[length=2mm]}, thick}
]
  \node (A) {A};
  \node (B) [right of=A] {B};
  \node (C) [below of=A] {C};
  \node (D) [right of=C] {D};
  \draw[edge] (A) -- node[above, font=\scriptsize] {4,2 km} (B);
  \draw[edge] (A) -- node[left, font=\scriptsize] {3,1 km} (C);
  \draw[edge] (B) -- node[right, font=\scriptsize] {2,5 km} (D);
  \draw[edge] (C) -- node[below, font=\scriptsize] {5,0 km} (D);
  \draw[edge] (A) -- node[above, sloped, font=\scriptsize] {6,4 km} (D);
\end{tikzpicture}
\caption{Ví dụ đơn giản về mô hình đồ thị có bốn địa điểm}
\label{fig:example-graph}
\end{figure}

\section{Trạng thái và quy tắc chuyển trạng thái}

\begin{itemize}[itemsep=3pt]
  \item \textbf{Đỉnh:} một địa điểm có mã định danh, tên, loại địa điểm và tọa
  độ địa lý.
  \item \textbf{Cạnh:} một đoạn đường cho phép di chuyển trực tiếp từ đỉnh
  nguồn đến đỉnh đích, kèm các thuộc tính dùng để tính chi phí.
  \item \textbf{Trạng thái đầu:} địa điểm xuất phát do người dùng chọn.
  \item \textbf{Trạng thái đích:} một địa điểm đích; với bài toán đa điểm, mục
  tiêu là ghé đủ tập địa điểm yêu cầu.
  \item \textbf{Phép chuyển:} từ đỉnh hiện tại, thuật toán chỉ được đi qua một
  cạnh đi ra chưa bị đóng. Quy tắc này tự nhiên bảo toàn hướng lưu thông.
\end{itemize}

\section{Thuộc tính của cạnh}
\label{sec:edge-attributes}

\begin{table}[h!]
\centering
\caption{Các thuộc tính chính của cạnh}
\begin{tabular}{p{4.0cm}p{9.2cm}}
\toprule
Thuộc tính & Ý nghĩa \\
\midrule
\texttt{distance\_km} & Chiều dài tuyến đường thực tế, đơn vị km. \\
\texttt{base\_time\_min} & Thời gian cơ sở suy ra từ chiều dài và vận tốc giả định, đơn vị phút. \\
\texttt{adjusted\_time\_min} & Thời gian sau khi áp dụng hệ số của kịch bản. \\
\texttt{congestion} & Mức ùn tắc hiệu dụng từ 1 đến 5. \\
\texttt{risk} & Tổng rủi ro cơ sở, mưa, sương mù và thi công. \\
\texttt{road\_type} & Các loại đường OpenStreetMap xuất hiện trên tuyến. \\
\texttt{route\_cost} & Chi phí chuẩn hóa theo profile tối ưu được chọn. \\
\bottomrule
\end{tabular}
\end{table}

\section{Hàm chi phí}

Với mỗi kịch bản, các đại lượng khoảng cách, thời gian hiệu dụng và rủi ro được
chuẩn hóa min--max về $[0,1]$ trên tập cạnh còn mở; ùn tắc được chia cho 5.
Chi phí của cạnh $e$ là

\begin{equation}
C(e)=\alpha d_n(e)+\beta t_n(e)+\gamma c_n(e)+\delta r_n(e),
\end{equation}

trong đó $d_n,t_n,c_n,r_n$ lần lượt là khoảng cách, thời gian, ùn tắc và rủi
ro đã chuẩn hóa. Các bộ trọng số trong chương trình được trình bày ở
Bảng~\ref{tab:optimization-profiles}.

\begin{table}[h!]
\centering
\caption{Trọng số của các profile tối ưu}
\label{tab:optimization-profiles}
\begin{tabular}{lcccc}
\toprule
Profile & $\alpha$ & $\beta$ & $\gamma$ & $\delta$ \\
\midrule
Shortest & 0,80 & 0,15 & 0,03 & 0,02 \\
Fastest & 0,10 & 0,75 & 0,10 & 0,05 \\
Balanced & 0,20 & 0,50 & 0,15 & 0,15 \\
Safest & 0,10 & 0,25 & 0,10 & 0,55 \\
\bottomrule
\end{tabular}
\end{table}

Các trọng số đều có tổng bằng 1. Profile Shortest nhấn mạnh khoảng cách,
Fastest nhấn mạnh thời gian, Safest ưu tiên rủi ro, còn Balanced phân bổ trọng
số cho cả bốn yếu tố. Tổng chi phí của một tuyến bằng tổng $C(e)$ của các cạnh
trên tuyến. Riêng lựa chọn khoảng cách hoặc thời gian trên giao diện có thể yêu
cầu thuật toán tối ưu trực tiếp theo \texttt{distance\_km} hoặc
\texttt{adjusted\_time\_min}.
